\section{Two model verifications}\label{sec:v7-instances}
\subsection{Markov-renewal orbit geometry}
Let $A$ be a primitive zero--one matrix on $b$ symbols, with Perron
root greater than one. Its legal one-sided sequences have a stationary
fully supported law $\mu$ such that every legal conditional child
probability is at least $\zeta>0$. Let $h_i$ map a compact set of diameter
at most one into pairwise $g$-separated subsets of itself. Assume word
scales $r_v$, $r_{\emptyw}=1$, satisfy
\begin{align}
 D_0^{-1}r_v|x-y|\le |h_vx-h_vy|&\le r_v|x-y|,\nonumber\\
 D_0^{-1}r_ur_v\le r_{uv}&\le r_ur_v,
 \qquad 0<r_-\le r_i\le r_+<1.                         \label{eq:v7-markov-metric}
\end{align}
The injective coding map identifies the legal sequences with a compact
set $K$, on which $T$ deletes the first symbol. Let $f:K\to\R^m$ obey
\begin{equation}
 |f(x)-f(y)|\le H|x-y|^\beta,\qquad
 \sum_{j=0}^J|f(T^jx)-f(T^jy)|^2\ge c_o|x-y|^{2\beta},
 \quad 0<\beta\le1.                                   \label{eq:v7-delay}
\end{equation}
The renewal age weights satisfy
$0<q_-\le w_{k+1}/w_k\le q_+<1$.
Set $\tau=2\beta$ and
\begin{equation}
 V(v)=\sum_{k<|v|}w_kr_{v_{k:}}^\tau+W(|v|),\qquad
 a_v=\mu[v]V(v),\qquad
 Z(x)=(\sqrt{w_k}f(T^kx))_{k\ge0}.                    \label{eq:v7-markov-energy}
\end{equation}

\begin{corollary}[Markov renewal as a structural instance]
\label{cor:v7-markov}
The data \eqref{eq:v7-markov-metric}--\eqref{eq:v7-markov-energy}
satisfy Theorem~\ref{thm:v7-tree} in its deletion orientation.
In particular $e_M^2(Z)\asymp R_M\asymp G_M$, with a stationary
finite-tree observer and an unrestricted randomized, nonstationary
converse. Forbidden edges are not completed, and unary vertices are
charged.
\end{corollary}
\begin{proof}
Every suffix scale is at most one and appending a letter decreases its
scale. If $|u|=m$ and $|v|=n$, monotonicity of $w$ gives
\begin{align*}
 V(v)-\left(\sum_{j<n}w_{m+j}r_{v_{j:}}^\tau+W(m+n)\right)
 &\ge r_v^\tau\sum_{k<m}w_k,\\
 V(uv)-\left(\sum_{j<n}w_{m+j}r_{v_{j:}}^\tau+W(m+n)\right)
 &\le r_+^\tau r_v^\tau\sum_{k<m}w_k.
\end{align*}
Therefore $V(uv)<V(v)$. Stationarity gives only
$\mu[uv]\le\mu[v]$, which is enough to give the strict energy
inequality. Equality of cylinder masses at a unique predecessor causes
no problem.

Appending a legal letter multiplies old spatial terms by a factor
between $D_0^{-\tau}r_-^\tau$ and $r_+^\tau$. The old tail $W(n)$
becomes $w_nr_i^\tau+W(n+1)$. Since
$w_n\ge(1-q_+)W(n)$ and $W(n+1)\ge q_-W(n)$, we have
\[
 \ell V(v)\le V(vi)\le\rho V(v),\quad
 \ell=\min(D_0^{-\tau}r_-^\tau,q_-),\quad
 \rho=1-(1-r_+^\tau)(1-q_+)<1.
\]
Use $a_-=\zeta\ell$ and sum the conditional masses to verify
\eqref{eq:v7-tree-child}. These inequalities also give monotonicity of
$D_v=\sqrt{V(v)}$. In a primitive graph with Perron root greater than
one, a deterministic path cannot run for $b$ repeated unary vertices:
a repeated vertex would generate a closed deterministic cycle. Thus
one can take $U=b$.

Introduce $d_\beta(x,y)^2=\sum_kw_k|T^kx-T^ky|^{2\beta}$.
H\"older continuity bounds the orbit distance from above by $H d_\beta$.
Summing the delayed lower bound \eqref{eq:v7-delay}, and using
$w_{k-j}\le q_-^{-j}w_k$, bounds $d_\beta$ from above by a constant
times the orbit distance. Points in one cylinder have squared
$d_\beta$ distance at most $V(v)$. Points whose longest common prefix
is $u$ have, before its end, distances at least $D_0^{-1}g r_{u_{k:}}$,
and at its end distance at least $g$. The last term controls $W(|u|)$
since $W(n)\le w_n/(1-q_+)$. This proves both image comparisons.
Finally the H\"older bound on each suffix and the bounded tail prove
\eqref{eq:v7-deletion-osc}. The structural theorem now supplies all
combinatorial, quantization, and causal conclusions.
\end{proof}

This corollary recovers the earlier Markov-renewal profile without
repeating its tree argument as a new model-specific proof. Gibbs
pressure identifies the profile exponent under the additional
thermodynamic assumptions of the previous Markov theorem, whose full
proof is retained in the companion. Pressure is a computation of this
instance, not an assumption of the foundational realization theorem.
For clarity, the passage from power bounds to an exponent uses only the
following elementary fact.

\begin{proposition}[Logarithmic squeeze]\label{prop:v7-exponent-squeeze}
If for every $\delta>0$ a positive sequence $R_M$ satisfies
$c_\delta M^{-\alpha-\delta}\le R_M\le
C_\delta M^{-\alpha+\delta}$ eventually, then
$-\log R_M/\log M\to\alpha$.
\end{proposition}
\begin{proof}
Taking logarithms gives lower and upper limits between
$\alpha-\delta$ and $\alpha+\delta$, because fixed logarithmic constants
divided by $\log M$ vanish. Let $\delta$ decrease to zero. In a
pressure-root proof one first chooses the auxiliary root parameter on
either side of its zero and only then lets $M$ tend to infinity; root
continuity then supplies the arbitrary $\delta$ used here.
\end{proof}

\subsection{An infinite hidden shift with noisy continual reports}
Let $(B_t)_{t\in\mathbb Z}$ be independent identically distributed Bernoulli variables. The hidden
state at time $t$ is the whole sequence
$X_t=(B_t,B_{t-1},\ldots)$. A binary memoryless channel reports $Y_t$
about $B_t$. Write
\[
 \gamma=\Pp(Y_t=1)\in(0,1),\qquad
 a_i=\Pp(B_t=1\mid Y_t=i),\qquad a_0\ne a_1.
\]
Fix $0<r<1$. The Hilbert-valued prediction target and its mean are
\begin{equation}
 U_t=(\sqrt{1-r^2}\,r^j B_{t-j})_{j\ge0},\qquad
 m_t=(\sqrt{1-r^2}\,r^j a_{Y_{t-j}})_{j\ge0}.           \label{eq:v7-shift-target}
\end{equation}
This squared loss can equally be implemented by drawing a fresh terminal
query $j$ with probability $(1-r^2)r^{2j}$ and scoring the prediction
of $B_{t-j}$. The independent query is a task index, not a new
observation of the hidden state. Its distribution and interface are
fixed in advance.

\begin{theorem}[Sharp register law for an infinite posterior]
\label{thm:v7-shift}
Put $p_0=1-\gamma$, $p_1=\gamma$ and let $s\in(0,1)$ be the unique
solution of
\begin{equation}
            r^{2s}(p_0^s+p_1^s)=1,
       \qquad \alpha=(1-s)/s.                           \label{eq:v7-shift-root}
\end{equation}
For the actual noisy shift experiment,
\begin{equation}
 R_M=B+\Theta\bigl((a_1-a_0)^2M^{-\alpha}\bigr),\qquad
 B=\sum_{i=0}^1p_i a_i(1-a_i).                          \label{eq:v7-shift-law}
\end{equation}
The constants in the excess comparison depend on $r,p_0,p_1$, not on
$a_0,a_1$. No posterior truncation is assumed in deriving this law.
For a fair hidden bit and a binary symmetric channel of error
$\varepsilon<1/2$, it becomes
\[
 R_M=\varepsilon(1-\varepsilon)
    +\Theta\bigl((1-2\varepsilon)^2M^{2\log r/\log2}\bigr).
\]
At $\varepsilon=1/2$ the predictive mean is constant and the excess is
zero with one state.
\end{theorem}
\begin{proof}
The full posterior is the product of Bernoulli laws with parameters
$a_{Y_{t-j}}$. Thus \eqref{eq:v7-shift-target} is the actual conditional
mean, and orthogonality gives the displayed $B$. Normalize its
nonconstant part by $a_1-a_0$. On the full binary past language take
$p_v=\prod p_{v_j}$, $D_v=r^{|v|}$ and $a_v=p_vr^{2|v|}$.
Children multiply energy by $p_i r^2$, their sum multiplies it by $r^2$,
and every nonempty prefix attachment strictly decreases it. If two
pasts first differ at position $n$, their normalized squared image
distance lies between $(1-r^2)r^{2n}$ and $r^{2n}$. All hypotheses of
Theorem~\ref{thm:v7-tree} hold in the insertion orientation.

For the exact order of $G_M$, set $q_i=(p_ir^2)^s$. Equation
\eqref{eq:v7-shift-root} says $q_0+q_1=1$. For every finite complete
cut, consistency of this Bernoulli measure gives $\sum_v a_v^s=1$.
By greedy balance, for a cut with $L$ leaves and maximum $\theta$,
$L(a_-\theta)^s\le1\le L\theta^s$.
Hence $G_M\asymp L\theta\asymp L^{-(1-s)/s}\asymp M^{-\alpha}$.
The root exists uniquely because the left side of
\eqref{eq:v7-shift-root} decreases continuously from two to $r^2<1$.
The fair case solves it explicitly, and the constant-mean case is
immediate. If only reports from time zero are supplied, the unobserved
past contributes at most $r^{2(t+1)}$ to the squared target norm at
time $t$. Its average vanishes; the finite-tree initialization in the
structural theorem still gives \eqref{eq:v7-shift-law}.
\end{proof}

This is not an exactly closed finite list of conditional means. Every
past report changes its own posterior coordinate, and arbitrarily many
coordinates can be queried. It also has no fixed-step common minorizer:
after $k$ hidden updates the old tail remains intact, and two different
initial tails give disjoint supports. A measure dominated by all those
kernels must be zero. As always, a bare Borel encoding of an infinite
sequence into one real number is possible; it is not a finite-cardinality
register and does not change the theorem.
